\documentclass[a4paper,12pt]{article}

\usepackage[applemac]{inputenc}
\usepackage[T1]{fontenc}
\usepackage[ngerman]{babel}
\usepackage[nospace,noadjust]{cite}
\usepackage{eurosym}
\usepackage{amssymb,amsmath,mathdots}
\usepackage{graphicx}
\usepackage{color}
\definecolor{kit}{cmyk}{1,0,0.6,0}
\usepackage{hyperref}
\usepackage{breqn}
\usepackage{changepage}
\hypersetup{
    pdftoolbar=true,
    pdfmenubar=true,
    pdfpagemode=UseOutlines,
    bookmarksnumbered=true,
    linktocpage=true,
    colorlinks=false,
    colorlinks=false
}

\setlength{\parindent}{0pt}
\parskip1.5ex

\begin{document}
\begin{titlepage}
    \begin{center}
        \vspace*{-80pt}
        \includegraphics[scale=0.25]{img/kit_logo.png}

        \vspace*{40pt}

        Fakult\text{\"a}t f\text{\"u}r Wirtschaftswissenschaften \\[1ex]
        Institut f\text{\"u}r Operations Research (IOR) \\[1ex]
        Optimierungsans\text{\"a}tze unter Unsicherheit \\[1ex]
        Prof. Dr. Steffen Rebennack          

        \vspace*{25pt}   

        \vspace*{35pt}

        \fboxsep 40pt
        \fboxrule 6pt
        \fcolorbox{kit}{white}{
            \parbox{80mm}{
                \begin{center}
                    \Large{Abgabge Rechner\text{\"u}bung}\\ 
                    \Large{Wintersemester 2024/25}
                \end{center}
            }
        }

        \vspace*{40pt}

        \normalsize

        Vorname Nachname\\
        Matr. Nr.:\\
        Studiengang (B.Sc.)\\[4ex]

        und \\[4ex]

        Vorname Nachname\\
        Matr. Nr.:\\
        Studiengang (B.Sc.)\\[4ex]
    \end{center}
\end{titlepage}

\newpage

\section*{L\"osungen zu Aufgabe 1}

\subsection*{Aufgabenteil a}

Der Quellcode ist in der Datei \href{../src/r1/Aufgabe1.gms}{Aufgabe1.gms} zu finden.

\subsection*{Aufgabenteil b}

Die entsprechende Ausf\"uhrung kann unter \href{../results/r1/Aufgabe1.lst}{Aufgabe1.lst} eingesehen werden. Daraus l\"asst sich entnehmen, dass der optimale Wert bei \EUR{5,169180} liegt und der optimale Punkt bei
\[
\begin{pmatrix}
\text{Apfel} \\ 
\text{Cornflakes} \\ 
\text{Karotten} \\ 
\text{Kartoffeln} \\ 
\text{K\"ase} \\ 
\text{Milch} \\ 
\text{Schokolade} \\ 
\text{Spinat} \\ 
\text{Steak}
\end{pmatrix}^T
=
\begin{pmatrix}
0 \\ 
3 \\ 
4 \\ 
0,4512 \\ 
0 \\ 
0 \\ 
2,0111 \\
0 \\
1
\end{pmatrix}^T
\]
zu finden ist.

\newpage

\section*{L\"osungen zu Aufgabe 2}

Die Unsicherheitsmenge l\"asst sich allgemein folgenderma\ss en beschreiben:

\[
\mathcal{U} = \left\{
\underbrace{
\left[
\begin{array}{c|c}
(c^0)^\top & d^0 \\ \hline
A^0 & b^0
\end{array}
\right]
}_{D^0}
+ \sum_{\ell=1}^L \zeta_\ell
\underbrace{
\left[
\begin{array}{c|c}
(c^\ell)^\top & d^\ell \\ \hline
A^\ell & b^\ell
\end{array}
\right]
}_{D^\ell}
\;\middle|\; \zeta \in \mathbb{Z}
\right\}
\]

F\"ur die Aufgabe werden unter anderem die Zielfunktion, gegeben durch
\[
    \min \quad 0.22x_1 + 0.18x_2 + 0.07x_3 + 0.14x_4 + 0.55x_5 + 0.1x_6 + 0.54x_7 + 0.28x_8 + 3.2x_9
\]
und die zweite Nebenbedingung 
\[
0.35x_1 + 7x_2 + x_3 + 2x_4 + 25x_5 + 3.5x_6 + 9x_7 + 2.5x_8 + 21x_9 \geq 56
\]
ben\"otigt.

Zuerst gilt es, die zweite Nebenbedingung in Standardform zu bringen. Dazu wird die Ungleichung mit $-1$ multipliziert.
\[
    -0.35x_1  -7x_2  -x_3  -2x_4 - 25x_5 - 3.5x_6 - 9x_7 - 2.5x_8 - 21x_9 \leq -56
\]

Aus den entsprechenden Vor\"uberlegungen l\"asst direkt die nominale Datenmatrix $D^0_2$ bilden:
\[
    D^0_2 = \begin{pmatrix}
        0.22 & 0.18 & 0.07 & 0.14 & 0.55 & 0.1 & 0.54 & 0.28 & 3.2 & 0 \\
        -0.35 & -7 & -1 & -2 & -25 & -3.5 & -9 & -2.5 & -21 & -56
    \end{pmatrix}
\]
        
Die Unsicherheiten betreffen s\"amtliche Koeffizienten der Zielfunktion, sowie die N\"ahrwerte f\"ur Proteine, Fette, Calcium und Vitamin B2. Zus\"atzlich gibt es eine Schwankung im Mindestbedarf f\"ur Proteine von $10g$. Alle Werte k\"onnen den Tabellen in der Aufgabenstellung entnommen werden.

Daher ist es entscheidend, die Unsicherheiten der Ziel- und Nebenfunktionen in absoluten Werten zu ermitteln. Die Unsicherheitsgr\"o\ss e kann anschlie\ss end wie folgt dargestellt werden:
\[
    \text{ZF} = \begin{pmatrix}
        0.06 & \underbrace{0.027}_{\mathrm{0.18 \times 0.15}} & \underbrace{0.014}_{\mathrm{0.07 \times 0.2}} & 0.04 & 0.1 & \underbrace{0.025}_{\mathrm{0.1 \times 0.25}} & \underbrace{0.216}_{\mathrm{0.54 \times 0.4}} & 0.1 & \underbrace{1.28}_{\mathrm{3.2 \times 0.4}}
    \end{pmatrix}        
\]
\[
    \text{NF} = \begin{pmatrix}
        \underbrace{-0.07}_{\mathrm{-0.35 \times 0.2}} & \underbrace{-0.7}_{\mathrm{-7 \times 0.1}} & \underbrace{-0.2}_{\mathrm{-1 \times 0.2}} & \underbrace{-0.1}_{\mathrm{-2 \times 0.05}} & \underbrace{-0.25}_{\mathrm{-25 \times 0.01}} & \underbrace{-0.35}_{\mathrm{-3.5 \times 0.1}} & \underbrace{-0.09}_{\mathrm{-9 \times 0.01}} & \underbrace{-0.25}_{\mathrm{-2.5 \times 0.1}} & \underbrace{-3.15}_{\mathrm{-21 \times 0.15}}
    \end{pmatrix}
\]

Ausgehend von diesen Werten lassen sich nun alle Shift-Matrizen $D^\ell$ berechnen. Die Anzahl der Unsicherheiten betr\"agt $L = 19$, woraus sich eine identische Anzahl an Shift-Matrizen ergibt. Diese sind wie folgt:

\[
D^1_2 = \begin{pmatrix}
    0.06 & 0 & 0 & 0 & 0 & 0 & 0 & 0 & 0 & 0 \\
    0 & 0 & 0 & 0 & 0 & 0 & 0 & 0 & 0 & 0
\end{pmatrix}
\]
\[
D^2_2 = \begin{pmatrix}
    0 & 0.027 & 0 & 0 & 0 & 0 & 0 & 0 & 0 & 0 \\
    0 & 0 & 0 & 0 & 0 & 0 & 0 & 0 & 0 & 0
\end{pmatrix}
\]
\[
D^3_2 = \begin{pmatrix}
    0 & 0 & 0.014 & 0 & 0 & 0 & 0 & 0 & 0 & 0 \\
    0 & 0 & 0 & 0 & 0 & 0 & 0 & 0 & 0 & 0
\end{pmatrix}
\]
\[
D^4_2 = \begin{pmatrix}
    0 & 0 & 0 & 0.04 & 0 & 0 & 0 & 0 & 0 & 0 \\
    0 & 0 & 0 & 0 & 0 & 0 & 0 & 0 & 0 & 0
\end{pmatrix}
\]
\[
D^5_2 = \begin{pmatrix}
    0 & 0 & 0 & 0 & 0.1 & 0 & 0 & 0 & 0 & 0 \\
    0 & 0 & 0 & 0 & 0 & 0 & 0 & 0 & 0 & 0
\end{pmatrix}
\]
\[
D^6_2 = \begin{pmatrix}
    0 & 0 & 0 & 0 & 0 & 0.025 & 0 & 0 & 0 & 0 \\
    0 & 0 & 0 & 0 & 0 & 0 & 0 & 0 & 0 & 0
\end{pmatrix}
\]
\[
D^7_2 = \begin{pmatrix}
    0 & 0 & 0 & 0 & 0 & 0 & 0.216 & 0 & 0 & 0 \\
    0 & 0 & 0 & 0 & 0 & 0 & 0 & 0 & 0 & 0
\end{pmatrix}
\]
\[
D^8_2 = \begin{pmatrix}
    0 & 0 & 0 & 0 & 0 & 0 & 0 & 0.1 & 0 & 0 \\
    0 & 0 & 0 & 0 & 0 & 0 & 0 & 0 & 0 & 0
\end{pmatrix}
\]
\[
D^9_2 = \begin{pmatrix}
    0 & 0 & 0 & 0 & 0 & 0 & 0 & 0 & 1.28 & 0 \\
    0 & 0 & 0 & 0 & 0 & 0 & 0 & 0 & 0 & 0
\end{pmatrix}
\]
\[
D^{10}_2 = \begin{pmatrix}
    0 & 0 & 0 & 0 & 0 & 0 & 0 & 0 & 0 & 0 \\
    0.07 & 0 & 0 & 0 & 0 & 0 & 0 & 0 & 0 & 0
\end{pmatrix}
\]
\[
D^{11}_2 = \begin{pmatrix}
    0 & 0 & 0 & 0 & 0 & 0 & 0 & 0 & 0 & 0 \\
    0 & 0.7 & 0 & 0 & 0 & 0 & 0 & 0 & 0 & 0
\end{pmatrix}
\]
\[
D^{12}_2 = \begin{pmatrix}
    0 & 0 & 0 & 0 & 0 & 0 & 0 & 0 & 0 & 0 \\
    0 & 0 & 0.2 & 0 & 0 & 0 & 0 & 0 & 0 & 0
\end{pmatrix}
\]
\[
D^{13}_2 = \begin{pmatrix}
    0 & 0 & 0 & 0 & 0 & 0 & 0 & 0 & 0 & 0 \\
    0 & 0 & 0 & 0.1 & 0 & 0 & 0 & 0 & 0 & 0
\end{pmatrix}
\]
\[
D^{14}_2 = \begin{pmatrix}
    0 & 0 & 0 & 0 & 0 & 0 & 0 & 0 & 0 & 0 \\
    0 & 0 & 0 & 0 & 0.25 & 0 & 0 & 0 & 0 & 0
\end{pmatrix}
\]
\[
D^{15}_2 = \begin{pmatrix}
    0 & 0 & 0 & 0 & 0 & 0 & 0 & 0 & 0 & 0 \\
    0 & 0 & 0 & 0 & 0 & 0.35 & 0 & 0 & 0 & 0
\end{pmatrix}
\]
\[
D^{16}_2 = \begin{pmatrix}
    0 & 0 & 0 & 0 & 0 & 0 & 0 & 0 & 0 & 0 \\
    0 & 0 & 0 & 0 & 0 & 0 & 0.09 & 0 & 0 & 0
\end{pmatrix}
\]
\[
D^{17}_2 = \begin{pmatrix}
    0 & 0 & 0 & 0 & 0 & 0 & 0 & 0 & 0 & 0 \\
    0 & 0 & 0 & 0 & 0 & 0 & 0 & 0.25 & 0 & 0
\end{pmatrix}
\]
\[
D^{18}_2 = \begin{pmatrix}
    0 & 0 & 0 & 0 & 0 & 0 & 0 & 0 & 0 & 0 \\
    0 & 0 & 0 & 0 & 0 & 0 & 0 & 0 & 3.15 & 0
\end{pmatrix}
\]
\[
D^{19}_2 = \begin{pmatrix}
    0 & 0 & 0 & 0 & 0 & 0 & 0 & 0 & 0 & 0 \\
    0 & 0 & 0 & 0 & 0 & 0 & 0 & 0 & 0 & 10
\end{pmatrix}
\]

\newpage

\section*{L\"osungen zu Aufgabe 3}

\subsection*{Aufgabenteil a - Intervall-Unsicherheit}

\subsubsection*{Umwandeln des Ursprungsproblems in Standardform}

Das Ursprungsproblem muss in folgende Struktur umgewandelt werden:
\[
(\text{LO}) \quad \left\{ \min_{x} \left\{ \mathbf{c}^\top \mathbf{x} + d \; \middle| \; A\mathbf{x} \leq \mathbf{b} \right\} \right\}
\]
Entsprechend l\"asst sich dieses dann nun folgenderma\ss en darstellen:
\[
\begin{aligned}
    \min & \quad 0.22x_1 + 0.18x_2 + 0.07x_3 + 0.14x_4 + 0.55x_5 + 0.16x_6 + 0.54x_7 + 0.28x_8 + 3.2x_9 \\
    \text{s.t.} & \quad -52x_1 - 355x_2 - 26x_3 - 71x_4 - 354x_5 - 64x_6 - 536x_7 - 17x_8 - 121x_9 \leq -2400 \\
    & \quad -0.35x_1 - 7x_2 - x_3 - 2x_4 - 25x_5 - 3.5x_6 - 9x_7 - 2.5x_8 - 21x_9 \leq -56 \\
    & \quad -18x_1 - 307x_2 - 7x_3 - 24.5x_4 + 177x_5 + 12x_6 + 52x_7 + 6.5x_8 + 60.5x_9 \leq 0 \\
    & \quad -0.4x_1 - 0.6x_2 - 0.2x_3 - 0.11x_4 - 28.3x_5 - 3.5x_6 - 31.5x_7 - 0.3x_8 - 4x_9 \leq -50 \\
    & \quad 0.4x_1 + 0.6x_2 + 0.2x_3 + 0.11x_4 + 28.3x_5 + 3.5x_6 + 31.5x_7 + 0.3x_8 + 4x_9 \leq 70 & \\
    & \quad -7x_1 - 13x_2 - 41x_3 - 6x_4 - 800x_5 - 120x_6 - 214x_7 - 126x_8 - 3x_9 \leq -500 \\
    & \quad -30x_1 - 60x_2 - 53x_3 - 47x_4 - 300x_5 - 170x_6 - 370x_7 - 230x_8 - 130x_9 \leq -1100 \\
    & \quad -x_1 - x_3 - x_8 \leq -4 \\
    & \quad x_1, \dots, x_9 \leq 5 \\
    & \quad - x_1, \dots, - x_9 \leq 0 \\
    & \quad x_2 \leq 3 \\
    & \quad x_6 \leq 2 \\
    & \quad x_7 \leq 3 \\
    & \quad - x_9 \leq - 1
\end{aligned}
\]

\subsubsection*{Bilden der nominalen Datenmatrix $D^0$}

Mit den vorherigen \"Uberlegungen l\"asst sich ganz einfach die nominale Datenmatrix $D^0$ dieser Struktur
\[
D^0 = \begin{pmatrix}\begin{array}{c|c}
(c^0)^\top & d^0 \\ \hline
A^0 & b^0
\end{array}\end{pmatrix}
\]
bilden.

Nach dem Streichen von Zeilen ohne Unsicherheiten lautet diese dann wie folgt:
\[
D^0 = \begin{pmatrix}
0.22 & 0.18 & 0.07 & 0.14 & 0.55 & 0.16 & 0.54 & 0.28 & 3.2 & 0 \\ % Zielfunktion
-52 & -355 & -26 & -71 & -354 & -64 & -536 & -17 & -121 & -2400 \\ % Kalorien
-0.35 & -7 & -1 & -2 & -25 & -3.5 & -9 & -2.5 & -21 & -56 \\ % Mind. Proteine
-0.4 & -0.6 & -0.2 & -0.11 & -28.3 & -3.5 & -31.5 & -0.3 & -4 & -50 \\ % Mind. Fett
0.4 & 0.6 & 0.2 & 0.11 & 28.3 & 3.5 & 31.5 & 0.3 & 4 & 70 \\ % Max. Fett
-7 & -13 & -41 & -6 & -800 & -120 & -214 & -126 & -3 & -500 \\ % Mind. Calcium
-30 & -60 & -53 & -47 & -300 & -170 & -370 & -230 & -130 & -1100 \\ % Mind. Vit. B2
\end{pmatrix}
\]

\subsubsection*{Bilden der Shiftmatrizen $D^\ell$}
Die Anzahl der Shiftmatrizen $D^\ell$ ergibt sich aus der Anzahl der Unsicherheiten. Alle Unsicherheiten lassen sich aus den Tabellen entnehmen.

F\"ur die Shiftmatrizen $D^\ell$ gilt es, die Unsicherheiten ggf. durch Multiplikation mit dem entsprechenden Koeffizienten der Zielfunktion in absolute Abweichungen umzuwandeln.
Wie in \"{U}bung $3f$ und Vorlesung $7$ (Abschnitt 5.3) erw\"{a}hnt, nutzt man f\"{u}r jede Unsicherheit in einer Zeile eine neue Shiftmatrix.
Durch diesen Ansatz lassen sich die Unsicherheiten eines Parameters $x_{\ell \in [1,9]}$ und des Mindestbedarfs $b$ in einer Shiftmatrix $D^\ell$ darstellen. Au\ss erdem ist es zu Vereinfachungszwecken erlaubt, Zeilen, in denen keine Unsicherheiten auftreten, zu streichen.

Die Shiftmatrizen lauten dann wie folgt:

\[
D^1_{x_1} = \begin{pmatrix}
0.06 & 0 & 0 & 0 & 0 & 0 & 0 & 0 & 0 & 0 \\ % Zielfunktion
0.07 & 0 & 0 & 0 & 0 & 0 & 0 & 0 & 0 & 0 \\ % Proteine
0.8 & 0 & 0 & 0 & 0 & 0 & 0 & 0 & 0 & 0 \\ % Mind. Fette
0.8 & 0 & 0 & 0 & 0 & 0 & 0 & 0 & 0 & 0 \\ % Max. Fette
1.4 & 0 & 0 & 0 & 0 & 0 & 0 & 0 & 0 & 0 \\ % Calcium
4.5 & 0 & 0 & 0 & 0 & 0 & 0 & 0 & 0 & 0 % Vit. B2
\end{pmatrix}
\begin{pmatrix}
    Zielfunktion \\
    Proteine \\
    Min. Fette \\
    Max. Fette \\
    Calcium \\
    Vit. B2
\end{pmatrix}
\]
\[
D^2_{x_2} = \begin{pmatrix}
0 & 0.027 & 0 & 0 & 0 & 0 & 0 & 0 & 0 & 0 \\ % Zielfunktion
0 & 0.7 & 0 & 0 & 0 & 0 & 0 & 0 & 0 & 0 \\ % Proteine
0 & 0.12 & 0 & 0 & 0 & 0 & 0 & 0 & 0 & 0 \\ % Mind. Fette
0 & 0.12 & 0 & 0 & 0 & 0 & 0 & 0 & 0 & 0 \\ % Max. Fette
0 & 6.5 & 0 & 0 & 0 & 0 & 0 & 0 & 0 & 0 \\ % Calcium
0 & 12 & 0 & 0 & 0 & 0 & 0 & 0 & 0 & 0 % Vit. B2
\end{pmatrix}
\begin{pmatrix}
    Zielfunktion \\
    Proteine \\
    Min. Fette \\
    Max. Fette \\
    Calcium \\
    Vit. B2
\end{pmatrix}
\]
\[
D^3_{x_3} = \begin{pmatrix} 
0 & 0 & 0.014 & 0 & 0 & 0 & 0 & 0 & 0 & 0 \\ % Zielfunktion
0 & 0 & 0.2 & 0 & 0 & 0 & 0 & 0 & 0 & 0 \\ % Proteine
0 & 0 & 0.04 & 0 & 0 & 0 & 0 & 0 & 0 & 0 \\ % Mind. Fette
0 & 0 & 0.04 & 0 & 0 & 0 & 0 & 0 & 0 & 0 \\ % Max. Fette
0 & 0 & 8.2 & 0 & 0 & 0 & 0 & 0 & 0 & 0 \\ % Calcium
0 & 0 & 10.6 & 0 & 0 & 0 & 0 & 0 & 0 & 0
\end{pmatrix}
\begin{pmatrix}
    Zielfunktion \\
    Proteine \\
    Min. Fette \\
    Max. Fette \\
    Calcium \\
    Vit. B2
\end{pmatrix}
\]
\[
D^4_{x_4} = \begin{pmatrix}
0 & 0 & 0 & 0.04 & 0 & 0 & 0 & 0 & 0 & 0 \\ % Zielfunktion
0 & 0 & 0 & 0.1 & 0 & 0 & 0 & 0 & 0 & 0 \\ % Proteine
0 & 0 & 0 & 0.0022 & 0 & 0 & 0 & 0 & 0 & 0 \\ % Mind. Fette
0 & 0 & 0 & 0.0022 & 0 & 0 & 0 & 0 & 0 & 0 \\ % Max. Fette
0 & 0 & 0 & 0.12 & 0 & 0 & 0 & 0 & 0 & 0 \\ % Calcium
0 & 0 & 0 & 0.47 & 0 & 0 & 0 & 0 & 0 & 0 % Vit. B2
\end{pmatrix}
\begin{pmatrix}
    Zielfunktion \\
    Proteine \\
    Min. Fette \\
    Max. Fette \\
    Calcium \\
    Vit. B2
\end{pmatrix}
\]
\[
D^5_{x_5} = \begin{pmatrix}
0 & 0 & 0 & 0 & 0.1 & 0 & 0 & 0 & 0 & 0 \\ % Zielfunktion
0 & 0 & 0 & 0 & 1.77 & 0 & 0 & 0 & 0 & 0 \\ % Proteine
0 & 0 & 0 & 0 & 0.035 & 0 & 0 & 0 & 0 & 0 \\ % Mind. Fette
0 & 0 & 0 & 0 & 0.035 & 0 & 0 & 0 & 0 & 0 \\ % Max. Fette
0 & 0 & 0 & 0 & 80 & 0 & 0 & 0 & 0 & 0 \\ % Calcium
0 & 0 & 0 & 0 & 15 & 0 & 0 & 0 & 0 & 0 % Vit. B2
\end{pmatrix}
\begin{pmatrix}
    Zielfunktion \\
    Proteine \\
    Min. Fette \\
    Max. Fette \\
    Calcium \\
    Vit. B2
\end{pmatrix}
\]
\[
D^6_{x_6} = \begin{pmatrix}
0 & 0 & 0 & 0 & 0 & 0.04 & 0 & 0 & 0 & 0 \\ % Zielfunktion
0 & 0 & 0 & 0 & 0 & 0.35 & 0 & 0 & 0 & 0 \\ % Proteine
0 & 0 & 0 & 0 & 0 & 0.35 & 0 & 0 & 0 & 0 \\ % Mind. Fette
0 & 0 & 0 & 0 & 0 & 0.35 & 0 & 0 & 0 & 0 \\ % Max. Fette
0 & 0 & 0 & 0 & 0 & 24 & 0 & 0 & 0 & 0 \\ % Calcium
0 & 0 & 0 & 0 & 0 & 8.5 & 0 & 0 & 0 & 0 % Vit. B2
\end{pmatrix}
\begin{pmatrix}
    Zielfunktion \\
    Proteine \\
    Min. Fette \\
    Max. Fette \\
    Calcium \\
    Vit. B2
\end{pmatrix}
\]
\[
D^7_{x_7} = \begin{pmatrix}
0 & 0 & 0 & 0 & 0 & 0 & 0.216 & 0 & 0 & 0 \\ % Zielfunktion
0 & 0 & 0 & 0 & 0 & 0 & 0.9 & 0 & 0 & 0 \\ % Proteine
0 & 0 & 0 & 0 & 0 & 0 & 3.15 & 0 & 0 & 0 \\ % Mind. Fette
0 & 0 & 0 & 0 & 0 & 0 & 3.15 & 0 & 0 & 0 \\ % Max. Fette
0 & 0 & 0 & 0 & 0 & 0 & 21.4 & 0 & 0 & 0 \\ % Calcium
0 & 0 & 0 & 0 & 0 & 0 & 37 & 0 & 0 & 0 % Vit. B2
\end{pmatrix}
\begin{pmatrix}
    Zielfunktion \\
    Proteine \\
    Min. Fette \\
    Max. Fette \\
    Calcium \\
    Vit. B2
\end{pmatrix}
\]
\[
D^8_{x_8} = \begin{pmatrix}
0 & 0 & 0 & 0 & 0 & 0 & 0 & 0.1 & 0 & 0 \\ % Zielfunktion
0 & 0 & 0 & 0 & 0 & 0 & 0 & 0.25 & 0 & 0 \\ % Proteine
0 & 0 & 0 & 0 & 0 & 0 & 0 & 0.045 & 0 & 0 \\ % Mind. Fette
0 & 0 & 0 & 0 & 0 & 0 & 0 & 0.045 & 0 & 0 \\ % Max. Fette
0 & 0 & 0 & 0 & 0 & 0 & 0 & 12.6 & 0 & 0 \\ % Calcium
0 & 0 & 0 & 0 & 0 & 0 & 0 & 46 & 0 & 0 % Vit. B2
\end{pmatrix}
\begin{pmatrix}
    Zielfunktion \\
    Proteine \\
    Min. Fette \\
    Max. Fette \\
    Calcium \\
    Vit. B2
\end{pmatrix}
\]
\[
D^9_{x_9} = \begin{pmatrix}
0 & 0 & 0 & 0 & 0 & 0 & 0 & 0 & 1.28 & 0 \\ % Zielfunktion
0 & 0 & 0 & 0 & 0 & 0 & 0 & 0 & 3.15 & 0 \\ % Proteine
0 & 0 & 0 & 0 & 0 & 0 & 0 & 0 & 1.2 & 0 \\ % Mind. Fette
0 & 0 & 0 & 0 & 0 & 0 & 0 & 0 & 1.2 & 0 \\ % Max. Fette
0 & 0 & 0 & 0 & 0 & 0 & 0 & 0 & 0.6 & 0 \\ % Calcium
0 & 0 & 0 & 0 & 0 & 0 & 0 & 0 & 19.5 & 0 % Vit. B2
\end{pmatrix}
\begin{pmatrix}
    Zielfunktion \\
    Proteine \\
    Min. Fette \\
    Max. Fette \\
    Calcium \\
    Vit. B2
\end{pmatrix}
\]
\[
D^{10}_{b} = \begin{pmatrix}
0 & 0 & 0 & 0 & 0 & 0 & 0 & 0 & 0 & 350 \\ % b_1
0 & 0 & 0 & 0 & 0 & 0 & 0 & 0 & 0 & 10 \\ % b_2
0 & 0 & 0 & 0 & 0 & 0 & 0 & 0 & 0 & 10 \\ % b_4
0 & 0 & 0 & 0 & 0 & 0 & 0 & 0 & 0 & 200 \\ % b_6
0 & 0 & 0 & 0 & 0 & 0 & 0 & 0 & 0 & 300 % b_7
\end{pmatrix}
\begin{pmatrix}
    Kalorien \\
    Proteine \\
    Min. Fette \\
    Calcium \\
    Vit. B2
\end{pmatrix}
\]

Insgesamt gilt dann:

\[
\mathcal{U} = \left\{ D^0 + \sum_{\ell=1}^{10}D^{\ell} \zeta_{\ell} \;\vert\; \zeta_{\ell} \in [-1, 1] \right\}
\]

\subsubsection*{Bilden des robusten Pendants}

Mit Hilfe der Unsicherheitsmenge aus dem vorherigen Schritt l\"asst sich nun das robuste Pendant des Ursprungsproblems bilden:
\[
\begin{aligned}
    &\min && \sup_{(c, d, \cdot) \in \mathcal{U}} \left(c^\top x + d\right) \\
    &\text{s.t.} && Ax \leq b \quad \forall (\cdot, A, b) \in \mathcal{U}, \\
    & && x \geq 0
\end{aligned}
\]
In diesem Fall ergibt es sich zu:
\[
\begin{aligned}
   \begin{aligned}
    \min & \quad \textcolor{blue}{c^\top x \quad \forall (c^\top, \cdot,\cdot,\cdot) \in \mathcal{U}} \\ % TODO
    \text{s.t.} & \quad \textcolor{blue}{a_1^\top x \leq b_1 \quad \forall (\cdot, \cdot, a_1^\top, b_1) \in \mathcal{U}} \\
    & \quad \textcolor{blue}{a_2^\top x \leq b_2 \quad \forall (\cdot, \cdot, a_2^\top, b_2) \in \mathcal{U}} \\
    & \quad -18x_1 - 307x_2 - 7x_3 - 24.5x_4 + 177x_5 + 12x_6 + 52x_7 + 6.5x_8 + 60.5x_9 \leq 0 \\
    & \quad \textcolor{blue}{a_4^\top x \leq b_4 \quad \forall (\cdot, \cdot, a_4^\top, b_4) \in \mathcal{U}} \\
    & \quad \textcolor{blue}{a_5^\top x \leq b_5 \quad \forall (\cdot, \cdot, a_5^\top, b_5) \in \mathcal{U}} \\
    & \quad \textcolor{blue}{a_6^\top x \leq b_6 \quad \forall (\cdot, \cdot, a_6^\top, b_6) \in \mathcal{U}} \\
    & \quad \textcolor{blue}{a_7^\top x \leq b_7 \quad \forall (\cdot, \cdot, a_7^\top, b_7) \in \mathcal{U}}\\
    & \quad -x_1 - x_3 - x_8 \leq -4 \\
    & \quad x_1, \dots, x_9 \leq 5 \\
    & \quad - x_1, \dots, - x_9 \leq 0 \\
    & \quad x_2 \leq 3 \\
    & \quad x_6 \leq 2 \\
    & \quad x_7 \leq 3 \\
    & \quad - x_9 \leq - 1
\end{aligned}
\end{aligned}
\]

\subsubsection*{Herleitung eines LP-\"Aquivalents}

Zur Veranschaulichung wird die zweite Nebenbedingung (Proteine) umgewandelt.
Dabei wird die Definition von $\mathcal{U}$ in die Nebenbedingung eingesetzt:
\[
(a_1^0)^\top x + \zeta_1 (a_1^1)^\top x + \dots + \zeta_9 (a_1^9)^\top x + \underbrace{\zeta_{10} (a_1^{10})^\top x}_{= 0} \leq b_1^0 + \underbrace{\zeta_1 b_1^1}_{=0} + \dots +\zeta_{10} b_1^{10} \quad \forall \zeta \in [-1, 1]
\]

Zuerst gilt es die Projektion der Unsicherheitsmenge bei Intervall-Unsicherheit zu bestimmen:
\[
\mathcal{U}^{2} = \left\{ 
\begin{array}{l}
    \begin{pmatrix}
    -0.35 \\ 
    -7 \\ 
    -1 \\
    -2 \\
    -25 \\
    -3.5 \\
    -9 \\
    -2.5 \\
    -21 \\
    -56
    \end{pmatrix}
    + \zeta_1 
    \begin{pmatrix}
    0.07 \\ 
    0 \\ 
    0 \\
    0 \\
    0 \\
    0 \\
    0 \\
    0 \\
    0 \\
    0
    \end{pmatrix}
    + \zeta_2 
    \begin{pmatrix}
    0 \\ 
    0.7 \\ 
    0 \\
    0 \\
    0 \\
    0 \\
    0 \\
    0 \\
    0 \\
    0
    \end{pmatrix}
    + \zeta_3 
    \begin{pmatrix}
    0 \\ 
    0 \\ 
    0.2 \\
    0 \\ 
    0 \\ 
    0 \\ 
    0 \\ 
    0 \\ 
    0 \\
    0
    \end{pmatrix}
    + \zeta_4 
    \begin{pmatrix}
    0 \\ 
    0 \\ 
    0 \\
    0.1 \\ 
    0 \\ 
    0 \\ 
    0 \\ 
    0 \\ 
    0 \\
    0
    \end{pmatrix}
    + \zeta_5 
    \begin{pmatrix}
    0 \\ 
    0 \\ 
    0\\
    0 \\ 
    1.77 \\ 
    0 \\ 
    0 \\ 
    0 \\ 
    0 \\
    0
    \end{pmatrix}
    \\[10pt]
    + \zeta_6 
    \begin{pmatrix}
    0 \\ 
    0 \\ 
    0 \\
    0 \\ 
    0 \\ 
    0.35 \\ 
    0 \\ 
    0 \\ 
    0 \\
    0
    \end{pmatrix}
    + \zeta_7 
    \begin{pmatrix}
    0 \\ 
    0 \\ 
    0 \\
    0 \\ 
    0 \\ 
    0 \\ 
    0.9 \\ 
    0 \\ 
    0 \\
    0
    \end{pmatrix}
    + \zeta_8 
    \begin{pmatrix}
    0 \\ 
    0 \\ 
    0 \\
    0 \\ 
    0 \\ 
    0 \\ 
    0 \\ 
    0.25 \\ 
    0 \\
    0
    \end{pmatrix}
    + \zeta_9
    \begin{pmatrix}
    0 \\ 
    0 \\ 
    0 \\
    0 \\ 
    0 \\ 
    0 \\ 
    0 \\ 
    0 \\ 
    3.15 \\
    0
    \end{pmatrix}
    + \zeta_{10} 
    \begin{pmatrix}
    0 \\ 
    0 \\ 
    0 \\
    0 \\ 
    0 \\ 
    0 \\ 
    0 \\ 
    0 \\ 
    0 \\
    10
    \end{pmatrix}
\end{array}
\; \middle| \; \|\zeta\|_\infty \leq 1 
\right\}.
\]

Damit ergibt sich f\"ur die Restriktion folgende robuste Darstellung:
\[
\begin{aligned}
    &\quad-0.35x_1 - 7x_2 - x_3 - 2x_4 - 25x_5 - 3.5x_6 - 9x_7 - 2.5x_8 - 21x_9 \\
    &\qquad + \max_{\|\zeta\|_\infty \leq 1} \big( 0.07\zeta_1 + 0.7\zeta_2 + 0.2\zeta_3 + 0.1\zeta_4 + 1.77\zeta_5 +  0.35\zeta_6\\
    &\qquad + 0.9\zeta_7 + 0.25\zeta_8 + 3.15\zeta_9 - 10\zeta_{10} \big) \leq -56
\end{aligned}
\]
\[
\begin{aligned}
    &\Leftrightarrow -0.35x_1 - 7x_2 - x_3 - 2x_4 - 25x_5 - 3.5x_6 - 9x_7 - 2.5x_8 - 21x_9 \\
    &\qquad + 0.07|\zeta_1| + 0.7|\zeta_2| + 0.2|\zeta_3| + 0.1|\zeta_4| + 1.77|\zeta_5| + 0.35|\zeta_6| \\
    &\qquad + 0.9|\zeta_7| + 0.25|\zeta_8| + 3.15|\zeta_9| + 10|\zeta_{10}| \leq -56
\end{aligned}
\]
\newpage
Ein lineares \"Aquivalent l\"asst sich daraufhin mit Hilfe des Lifting-Ansatzes aufstellen:
\[
\begin{aligned}
    &-0.35x_1 - 7x_2 - x_3 - 2x_4 - 25x_5 - 3.5x_6 - 9x_7 - 2.5x_8 - 21x_9 \\
    &\quad + w_{21} + w_{22} + w_{23} + w_{24} + w_{25} + w_{26} + w_{27}+ w_{28}+ w_{29}+ w_{210}
\end{aligned}
\begin{aligned}
    &\quad \leq -56
\end{aligned}
\]
\[
\begin{aligned} 
    -w_{21} \leq 0.07x_1 \leq w_{21} \\
    -w_{22} \leq 0.7x_2 \leq w_{22} \\
    -w_{23} \leq 0.2x_3 \leq w_{23} \\
    -w_{24} \leq 0.1x_4 \leq w_{24} \\
    -w_{25} \leq 1.77x_5 \leq w_{25} \\
    -w_{26} \leq 0.35x_6 \leq w_{26} \\
    -w_{27} \leq 0.9x_7 \leq w_{27} \\
    -w_{28} \leq 0.25x_8 \leq w_{28} \\
    -w_{29} \leq 3.15x_9 \leq w_{29} \\
    -w_{210} \leq 10 \leq w_{210} \\
\end{aligned}
\]

Die restlichen Schritte zur Herleitung des linearen \"Aquivalents wurden analog durchgef\"uhrt. Man findet das gesamte lineare Programm unter dem Abschnitt \hyperref[sec:lp-equivalent-interval]{LP-\"Aquivalent bei Intervall-Unsicherheit}.

\newpage

\subsubsection*{L\"osung des linearen Programms}

Nach der Herleitung des linearen \"Aquivalents kann das lineare Programm gel\"ost werden. Hierzu wurde der Code entsprechend in \href{../src/r3/Aufgabe3a.gms}{Aufgabe3a.gms} umgesetzt.
Jedoch muss erw\"ahnt werden, dass es uns nicht m\"oglich war, erfolgreich das lineare Programm in GAMS umzusetzen. Dies merkten wir daran, dass der optimale Wert entgegen unserer Intuition bei \EUR{13.610} lag. Aus diesem Grund haben wir uns dazu entschieden, das entsprechende Problem zus\"atzlich mit dem \href{../src/r3/Aufgabe3a_NaiveApproach.gms}{naiven Ansatz} zu l\"osen, auf welchem die nachfolgenden Erkenntnisse basieren.
Daraus l\"asst sich entnehmen, dass die optimale L\"osung in folgendem Punkt liegt:
\[
x^* = \begin{pmatrix}
    \text{Apfel} \\ \text{Cornflakes} \\ \text{Karotten} \\ \text{Kartoffeln} \\ \text{K\"ase} \\ \text{Milch} \\ \text{Schokolade} \\ \text{Spinat} \\ \text{Steak}
    \end{pmatrix}^T
     =
     \begin{pmatrix}
     0, 3, 4, 4.7749, 0, 1.0713, 1.9635, 0, 1
        \end{pmatrix}
\]
Der optimale Zielfunktionswert betr\"agt \EUR{7.914778}.

\subsubsection*{Zusammensetzung der Variablen und Restriktionen} 

Die Variablen und Restriktionen des linearen Programms besteht aus $9$ Entscheidungsvariablen $x_i$, $51$ Unsicherheitsvariablen in den Nebenbedingungen $w_{ai}$ und $9$ Unsicherheitsvariablen in der Zielfunktion $w_{0i}$.
Zu den vorher bestehenden Restriktionen kommen je Unsicherheitsvariable eine Restriktion hinzu, die die Unsicherheiten beschreiben, d.h. es kommen auf die bestehenden $14$ Restriktionen $60$ hinzu.
Demnach besteht das lineare Programm aus $69$ Variablen und $74$ Restriktionen.

\newpage

\subsection*{Aufgabenteil b - Ellips-Unsicherheit}

\subsubsection*{Herleitung eines LP-\"Aquivalents}

Um das lineare \"Aquivalent f\"ur Ellips-Unsicherheit abzuleiten, k\"onnen wir \"ahnliche Vor\"uberlegungen wie bei Intervall-Unsicherheit anwenden. Allerdings wird die Unsicherheitsmenge $\mathcal{U}$ unter Verwendung der euklidischen Norm anstelle der Supremumsnorm beschrieben. In diesem Fall lautet sie dann wie folgt:
\[
\mathcal{U} = \left\{ D^0 + \sum_{\ell=1}^{10}D^{\ell} \zeta_{\ell} \;\vert\; \|\zeta\|_2 \leq 1 \right\}
\]
Wir betrachten dieselbe Nebenbedingung wie im Fall der Intervall-Unsicherheit und erhalten das folgende robuste Pendant:
\[
\begin{aligned}
    &\quad -0.35x_1 - 7x_2 - x_3 - 2x_4 - 25x_5 - 3.5x_6 - 9x_7 - 2.5x_8 - 21x_9 \\
    &\quad + \max_{\|\zeta\|_2 \leq 1} \big( 0.07\zeta_1 + 0.7\zeta_2 + 0.2\zeta_3 + 0.1\zeta_4 + 1.77\zeta_5 +  0.35\zeta_6\\
    &\qquad + 0.9\zeta_7 + 0.25\zeta_8 + 3.15\zeta_9 - 10\zeta_{10} \big) \leq -56
\end{aligned}
\]
Da wir es mit der euklidischen Norm zu tun haben, kann die Nebenbedingung wie folgt umformuliert werden:
\[
\begin{aligned}
    &\quad -0.35x_1 - 7x_2 - x_3 - 2x_4 - 25x_5 - 3.5x_6 - 9x_7 - 2.5x_8 - 21x_9 \\
    &\qquad + \sqrt{
        \begin{aligned}0.0049x_1^2 + 0.49x_2 + 0.04x_3^2 + 0.01x_4^2 + 3.1329x_5^2 \\
            + 0.1225x_6^2 + 0.81x_7^2 + 0.0625x_8^2 + 9.9225x_9^2 + 100 \end{aligned}
            } \leq -56
\end{aligned}
\]

Die restlichen Schritte zur Herleitung des linearen \"Aquivalents wurden analog durchgef\"uhrt. Man findet das gesamte lineare Programm unter dem Abschnitt \hyperref[sec:lp-equivalent-ellips]{LP-\"Aquivalent bei Ellips-Unsicherheit}.

\subsubsection*{L\"osung des Programms}

Nach der Herleitung des \"Aquivalents kann das Programm gel\"ost werden. Hierzu wurde der Code entsprechend in \href{../src/r3/Aufgabe3b.gms}{Aufgabe3b.gms} umgesetzt. Da es sich um ein nichtlineares Problem handelt, wird es mit einem entsprechenden Solver gel\"ost.
Aus der Durchf\"uhrung l\"asst sich entnehmen, dass die optimale L\"osung in folgendem Punkt liegt:
\[
x^* = \begin{pmatrix}
    \text{Apfel} \\ \text{Cornflakes} \\ \text{Karotten} \\ \text{Kartoffeln} \\ \text{K\"ase} \\ \text{Milch} \\ \text{Schokolade} \\ \text{Spinat} \\ \text{Steak}
    \end{pmatrix}^T
     =
     \begin{pmatrix}
     2.484, 3, 5, 5, 0, 0.371, 1.728, 0, 1
        \end{pmatrix}
\]
Der optimale Zielfunktionswert betr\"agt \EUR{7.6895}.

\subsection*{Aufgabenteil c - Vergleich der Ergebnisse}

Wie aus den vorherigen Teilaufgaben ersichtlich, unterscheiden sich die Ergebnisse der verschiedenen Unsicherheitsmodelle:
\begin{itemize}
    \item Nominalfall ohne Unsicherheiten: \EUR{5.1692}
    \item Intervall-Unsicherheit: \EUR{7.9148}
    \item Ellips-Unsicherheit: \EUR{7.6895}
\end{itemize}
Auf den ersten Blick erkennt man, dass die Ber\"ucksichtigung von Unsicherheiten die Zielwerte erh\"oht, da robuste Optimierungen mehr Sicherheit bieten, jedoch auf Kosten der Effizienz.
\begin{description}
    \item[1b) und 3a):] Hier besteht ein deutlicher Unterschied, der zeigt, wie stark die Annahme von Intervall-Unsicherheiten den Zielwert beeinflusst. Der Wert steigt um ca. $53 \%$, da die Worst-Case-Kombinationen der Parameter einbezogen werden.
    \item[1b) und 3b):] Auch hier steigt der Zielwert deutlich, jedoch moderater als bei der Intervall-Unsicherheit, um ca. $49 \%$. 
    \item[3a) und 3b):] Der Unterschied zwischen Intervall- und Ellips-Unsicherheit betr\"agt ca. $2.85\%$. Man erkennt, dass die Ellips-Unsicherheit effizienter ist, da sie realistischere Abh\"angigkeiten und Verteilungen der Unsicherheiten ber\"ucksichtigt.
\end{description}
Die Unterschiede in den Ergebnissen zeigen, dass die Wahl des Unsicherheitsmodells einen gro\ss en Einfluss auf die Zielfunktion hat. Interpretieren l\"asst sich dies folgenderma\ss en: 
\begin{description}
    \item[Keine Unsicherheit:] Die L\"osung ohne Unsicherheiten ist die g\"unstigste, jedoch nicht robust. Sobald Unsicherheiten auftreten, kann die L\"osung unbrauchbar werden.
    \item[Intervall-Unsicherheit:] Die L\"osung mit Intervall-Unsicherheiten ist sehr konservativ. Sie ber\"ucksichtigt Worst-Case-Szenarien, bei denen alle Unsicherheiten ihre Extremwerte annehmen k\"onnen. Dadurch steigt der Zielwert stark an, was hohe Sicherheit, aber auch Ineffizienz bedeutet.
    \item[Ellips-Unsicherheit:] Diese L\"osung bietet ebenfalls Robustheit, ist jedoch weniger konservativ als die L\"osung mit Intervall-Unsicherheiten. Die Ber\"ucksichtigung von Abh\"angigkeiten zwischen Unsicherheiten f\"uhrt zu einem niedrigeren Zielwert, was im Kontext der Minimierung besser ist.
\end{description}
Abschlie\ss end l\"asst sich sagen, dass das Einbeziehen von Unsicherheiten die Zielwerte erheblich erh\"oht, wobei die Ber\"ucksichtigung von Intervall Unsicherheiten die konservativste, aber auch teuerste L\"osung bietet. Der Unterschied zwischen den robusten L\"osungen zeigt, dass die Wahl der Unsicherheitsmenge einen Einfluss auf den Zielwert hat. Dieser Unterschied ist jedoch nicht extrem und eher von Relevanz, wenn Effizienz eine besonders wichtige Rolle spielt.

\subsection*{Aufgabenteil d - Erh\"ohung der ben\"otigten Kalorien}

Nach Einf\"uhrung der Intervall- und Ellips-Unsicherheiten l\"asst sich der erh\"ohte Kalorienbedarf nicht mehr decken. Das liegt daran, dass durch die erh\"ohte Kalorienzufuhr die maximale Fett-Restriktion verletzt wird, da sie eine obere Schranke darstellt. Grund daf\"ur ist, dass das robuste Pendant Worst-Case-Szenarien abbildet. In einem Szenario ohne Unsicherheiten k\"onnte der erh\"ohte Kalorienbedarf gedeckt werden.
Dies l\"asst sich durch Anpassen der Restriktionen in den bisherigen GAMS-Dateien \"uberpr\"ufen.

\newpage

\subsubsection*{LP-\"Aquivalent bei Intervall-Unsicherheit}
\label{sec:lp-equivalent-interval}
\begin{adjustwidth}{-1.5cm}{}
\begin{tiny}
\setlength{\jot}{0pt}
\[
\begin{aligned}
   \begin{aligned}
    \min & \quad  0.22x_1 + 0.18x_2 + 0.07x_3 + 0.14x_4 + 0.55x_5 + 0.16x_6 + 0.54x_7 + 0.28x_8 + 3.2x_9 + w_{01} + w_{02} + w_{03} + w_{04} + w_{05} + w_{06} + w_{07}+ w_{08}+ w_{09} \\
    \text{s.t.} & \quad -w_{01} \leq 0.06x_1 \leq w_{01} \\
    & \quad -w_{02} \leq 0.027x_2 \leq w_{02} \\
    & \quad -w_{03} \leq 0.014x_3 \leq w_{03} \\
    & \quad -w_{04} \leq 0.04x_4 \leq w_{04} \\
    & \quad -w_{05} \leq 0.1x_5 \leq w_{05} \\
    & \quad -w_{06} \leq 0.04x_6 \leq w_{06} \\
    & \quad -w_{07} \leq 0.216x_7 \leq w_{07} \\
    & \quad -w_{08} \leq 0.1x_8 \leq w_{08} \\
    & \quad -w_{09} \leq 1.28x_9 \leq w_{09} \\
    & \quad -52x_1-355x_2-26x_3-71x_4-354x_5-64x_6-536x_7-17x_8-121x_9 \leq -2400 - w_{110}\\
    & \quad -w_{110} \leq 350 \leq w_{110} \\
    & \quad -0.35x_1 - 7x_2 - x_3 - 2x_4 - 25x_5 - 3.5x_6 - 9x_7 - 2.5x_8 - 21x_9 + w_{21} + w_{22} + w_{23} + w_{24} + w_{25} + w_{26} + w_{27}+ w_{28}+ w_{29} \leq -56 - w_{210} \\
    & \quad -w_{21} \leq 0.07x_1 \leq w_{21} \\
    & \quad -w_{22} \leq 0.7x_2 \leq w_{22} \\
    & \quad -w_{23} \leq 0.2x_3 \leq w_{23} \\
    & \quad -w_{24} \leq 0.1x_4 \leq w_{24} \\
    & \quad -w_{25} \leq 1.77x_5 \leq w_{25} \\
    & \quad -w_{26} \leq 0.35x_6 \leq w_{26} \\
    & \quad -w_{27} \leq 0.9x_7 \leq w_{27} \\
    & \quad -w_{28} \leq 0.25x_8 \leq w_{28} \\
    & \quad -w_{29} \leq 3.15x_9 \leq w_{29} \\
    & \quad -w_{210} \leq 10 \leq w_{210} \\
    & \quad -18x_1 - 307x_2 - 7x_3 - 24.5x_4 + 177x_5 + 12x_6 + 52x_7 + 6.5x_8 + 60.5x_9 \leq 0 \\
    & \quad -0.4x_1 - 0.6x_2 - 0.2x_3 - 0.11x_4 - 28.3x_5 - 3.5x_6 - 31.5x_7 - 0.3x_8 - 4x_9  + w_{41} + w_{42} + w_{43} + w_{44} + w_{45} + w_{46} + w_{47}+ w_{48}+ w_{49} \leq -50 - w_{410} \\
    & \quad -w_{41} \leq 0.8x_1 \leq w_{41} \\
    & \quad -w_{42} \leq 0.12x_2 \leq w_{42} \\
    & \quad -w_{43} \leq 0.04x_3 \leq w_{43} \\
    & \quad -w_{44} \leq 0.0022x_4 \leq w_{44} \\
    & \quad -w_{45} \leq 0.035x_5 \leq w_{45} \\
    & \quad -w_{46} \leq 0.35x_6 \leq w_{46} \\
    & \quad -w_{47} \leq 3.15x_7 \leq w_{47} \\
    & \quad -w_{48} \leq 0.045x_8 \leq w_{48} \\
    & \quad -w_{49} \leq 1.2x_9 \leq w_{49} \\
    & \quad -w_{410} \leq 10 \leq w_{410} \\
    & \quad 0.4x_1 + 0.6x_2 + 0.2x_3 + 0.11x_4 + 28.3x_5 + 3.5x_6 + 31.5x_7 + 0.3x_8 + 4x_9 + w_{51} + w_{52} + w_{53} + w_{54} + w_{55} + w_{56} + w_{57}+ w_{58}+ w_{59} \leq 70 \\
    & \quad -w_{51} \leq 0.8x_1 \leq w_{51} \\
    & \quad -w_{52} \leq 0.12x_2 \leq w_{52} \\
    & \quad -w_{53} \leq 0.04x_3 \leq w_{53} \\
    & \quad -w_{54} \leq 0.0022x_4 \leq w_{54} \\
    & \quad -w_{55} \leq 0.035x_5 \leq w_{55} \\
    & \quad -w_{56} \leq 0.35x_6 \leq w_{56} \\
    & \quad -w_{57} \leq 3.15x_7 \leq w_{57} \\
    & \quad -w_{58} \leq 0.045x_8 \leq w_{58} \\
    & \quad -w_{59} \leq 1.2x_9 \leq w_{59} \\
    & \quad -7x_1 - 13x_2 - 41x_3 - 6x_4 - 800x_5 - 120x_6 - 214x_7 - 126x_8 - 3x_9 + w_{61} + w_{62} + w_{63} + w_{64} + w_{65} + w_{66} + w_{67}+ w_{68}+ w_{69} \leq -500 - w_{610} \\
    & \quad -w_{61} \leq 1.4x_1 \leq w_{61} \\
    & \quad -w_{62} \leq 6.5x_2 \leq w_{62} \\
    & \quad -w_{63} \leq 8.2x_5 \leq w_{63} \\
    & \quad -w_{64} \leq 0.12x_4 \leq w_{64} \\
    & \quad -w_{65} \leq 80x_5 \leq w_{65} \\
    & \quad -w_{66} \leq 24x_6 \leq w_{66} \\
    & \quad -w_{67} \leq 21.4x_7 \leq w_{67} \\
    & \quad -w_{68} \leq 12.6x_8 \leq w_{68} \\
    & \quad -w_{69} \leq 0.6x_9 \leq w_{69} \\
    & \quad -w_{610} \leq 200 \leq w_{610} \\
    & \quad -30x_1 - 60x_2 - 53x_3 - 47x_4 - 300x_5 -170x_6 - 370x_7 -230x_8 - 130x_9 + w_{71} + w_{72} + w_{73} + w_{74} + w_{75} + w_{76} + w_{77}+ w_{78}+ w_{79} \leq -1100 - w_{710} \\
    & \quad -w_{71} \leq 4.5x_1 \leq w_{71} \\
    & \quad -w_{72} \leq 12x_2 \leq w_{72} \\
    & \quad -w_{73} \leq 10.6x_5 \leq w_{73} \\
    & \quad -w_{74} \leq 0.47x_4 \leq w_{74} \\
    & \quad -w_{75} \leq 15x_5 \leq w_{75} \\
    & \quad -w_{76} \leq 8.5x_6 \leq w_{76} \\
    & \quad -w_{77} \leq 37x_7 \leq w_{77} \\
    & \quad -w_{78} \leq 46x_8 \leq w_{78} \\
    & \quad -w_{79} \leq 19.5x_9 \leq w_{79} \\
    & \quad -w_{710} \leq 300 \leq w_{710} \\
    & \quad -x_1 - x_3 - x_8 \leq -4 \\
    & \quad x_1, \dots, x_9 \leq 5 \\
    & \quad - x_1, \dots, - x_9 \leq 0 \\
    & \quad x_2 \leq 3 \\
    & \quad x_6 \leq 2 \\
    & \quad x_7 \leq 3 \\
    & \quad - x_9 \leq - 1
    \end{aligned}
\end{aligned}
\]
\end{tiny}
\end{adjustwidth}

\newpage

\subsubsection*{LP-\"Aquivalent bei Ellips-Unsicherheit}
\label{sec:lp-equivalent-ellips}
\begin{adjustwidth}{-1.5cm}{}
\begin{tiny}
\setlength{\jot}{0pt}
\[
\begin{aligned}
   \begin{aligned}
    \min & \quad  0.22x_1 + 0.18x_2 + 0.07x_3 + 0.14x_4 + 0.55x_5 + 0.16x_6 + 0.54x_7 + 0.28x_8 + 3.2x_9 \\
    & \qquad+ \sqrt{0.0036x_1^2+0.000729x_2^2+0.000196x_3^2+0.0016x_4^2+0.01x_5^2+0.0016x_6^2+0.046656x_7^2+0.01x_8^2+1.6384x_9^2} \\
    \text{s.t.}
    & \quad -52x_1-355x_2-26x_3-71x_4-354x_5-64x_6-536x_7-17x_8-121x_9+\sqrt{122500} \leq -2400 \\
    & \quad -0.35x_1 - 7x_2 - x_3 - 2x_4 - 25x_5 - 3.5x_6 - 9x_7 - 2.5x_8 - 21x_9 \\
    & \qquad + \sqrt{0.0049x_1^2 + 0.49x_2^2 + 0.04x_3^2 + 0.01x_4^2 + 3.1329x_5^2 + 0.1225x_6^2 + 0.81x_7^2 + 0.0625x_8^2 + 9.9225x_9^2 + 100} \leq -56 \\
    & \quad -18x_1 - 307x_2 - 7x_3 - 24.5x_4 + 177x_5 + 12x_6 + 52x_7 + 6.5x_8 + 60.5x_9 \leq 0 \\
    & \quad -0.4x_1 - 0.6x_2 - 0.2x_3 - 0.11x_4 - 28.3x_5 - 3.5x_6 - 31.5x_7 - 0.3x_8 - 4x_9 \\
    & \qquad + \sqrt{0.64x_1^2+ 0.0144x_2^2 + 0.0016x_3^2 + 0.00000484x_4^2 + 0.001225x_5^2 + 0.1225x_6^2 + 9.9225x_7^2 + 0.002025x_8^2 + 1.44x_9^2 + 100} \leq -50 \\
    & \quad 0.4x_1 + 0.6x_2 + 0.2x_3 + 0.11x_4 + 28.3x_5 + 3.5x_6 + 31.5x_7 + 0.3x_8 + 4x_9 \\
    & \qquad + \sqrt{0.64x_1^2+ 0.0144x_2^2 + 0.0016x_3^2 + 0.00000484x_4^2 + 0.001225x_5^2 + 0.1225x_6^2 + 9.9225x_7^2 + 0.002025x_8^2 + 1.44x_9^2} \leq 70 \\
    & \quad -7x_1 - 13x_2 - 41x_3 - 6x_4 - 800x_5 - 120x_6 - 214x_7 - 126x_8 - 3x_9 \\
    & \qquad + \sqrt{1.96x_1^2 + 42.25x_2^2 + 67.24x_3^2 + 0.0144x_4^2 + 6400x_5^2 + 576x_6^2 + 457.96x_7^2 + 158.76x_8^2 + 0.36x_9^2 + 40000} \leq -500 \\
    & \quad -30x_1 - 60x_2 - 53x_3 - 47x_4 - 300x_5 -170x_6 - 370x_7 -230x_8 - 130x_9 \\
    & \qquad + \sqrt{20.25x_1^2 + 144x_2^2 + 112.36x_5^2 + 0.2209x_4^2 + 225x_5^2 + 72.25x_6^2 + 1369x_7^2 + 2116x_8^2 + 380.25x_9^2 + 90000} \leq -1100 \\
    & \quad -x_1 - x_3 - x_8 \leq -4 \\
    & \quad x_1, \dots, x_9 \leq 5 \\
    & \quad - x_1, \dots, - x_9 \leq 0 \\
    & \quad x_2 \leq 3 \\
    & \quad x_6 \leq 2 \\
    & \quad x_7 \leq 3 \\
    & \quad - x_9 \leq - 1
    \end{aligned}
\end{aligned}
\]
\end{tiny}
\end{adjustwidth}

\end{document}